\section*{अध्याय \ref{ch:g7:circuits-series-parallel} --- परिपथ: फंदे, श्रेणी, समांतर}

\begin{solution}{exo:g7:circuits-series-parallel:1}
\omterm{def:g3:first-electric-circuit:battery}{बैटरी} के धक्के से परिपथ में क़तार बाँधकर चलते आवेशित कणों का बहाव। वह
तभी तक रहती है जब तक कम से कम एक बंद फंदा दोनों सिरों को जोड़े रखे;
फंदा खोलो और क़तार एक साथ हर जगह रुक जाती है।
\end{solution}

\begin{solution}{exo:g7:circuits-series-parallel:2}
\omterm{def:g3:first-electric-circuit:battery}{बैटरी} के बाहर, $+$ से परिपथ का चक्कर लगाकर $-$ तक। एक \omterm{def:g3:first-electric-circuit:bulb}{बल्ब} वाले फंदे
में तीर $+$ से निकलते हैं, \omterm{def:g3:first-electric-circuit:bulb}{बल्ब} से गुज़रते हैं, और $-$ पर लौट आते
हैं।
\end{solution}

\begin{solution}{exo:g7:circuits-series-parallel:3}
इस ग़लती का, कि उपकरण धारा खा जाते हैं और इसलिए ``आगे वाले'' \omterm{def:g3:first-electric-circuit:bulb}{बल्ब} मंद
चमकेंगे। बराबर चमक --- जो हर अदला-बदली के बाद भी बनी रहती है ---
दिखाती है कि दोनों में से वही पूरी धारा गुज़रती है: रास्ते में कुछ भी
नहीं खाया जाता।
\end{solution}

\begin{solution}{exo:g7:circuits-series-parallel:4}
वे एक ही धारा हैं: अकेला फंदा अपने हर उपकरण में से एक ही क़तार ले जाता
है।
\end{solution}

\begin{solution}{exo:g7:circuits-series-parallel:5}
वह बँट जाती है --- हर \omterm{def:g5:series-parallel-circuits:parallel}{शाखा} क़तार का एक हिस्सा लेती है --- और दूसरी
\omterm{def:g5:series-parallel-circuits:parallel}{संधि} पर सारे हिस्से मिलकर फिर पूरे हो जाते हैं। हर \omterm{def:g5:series-parallel-circuits:parallel}{शाखा} को \omterm{def:g3:first-electric-circuit:battery}{बैटरी} का
पूरा धक्का मिलता है।
\end{solution}

\begin{solution}{exo:g7:circuits-series-parallel:6}
अकेला $=$ एक समांतर पड़ोसी के साथ (पूरी चमक), फिर एक श्रेणी वाले साथी
के साथ (मंद), और फिर दो श्रेणी वाले साथियों के साथ (और भी मंद)।
\end{solution}

\begin{solution}{exo:g7:circuits-series-parallel:7}
समांतर: दो शाखाएँ, और हर एक पूरी क़तार खींचती हुई, \omterm{def:g3:first-electric-circuit:battery}{बैटरी} की कुल आपूर्ति
दोगुनी कर देती हैं, जबकि श्रेणी वाली जोड़ी दोनों में से एक ही गला हुआ
क़तार भेजती है। पूरी चमक का दाम \omterm{def:g3:first-electric-circuit:battery}{बैटरी} पर चुकाया जाता है।
\end{solution}

\begin{solution}{exo:g7:circuits-series-parallel:8}
कुर्सियाँ आवेशित चलने वाले हैं, जो लगातार चक्कर लगाते हैं और कभी ख़र्च
नहीं होते; स्की करने वाले \omterm{def:g5:everyday-energy:energy}{ऊर्जा} हैं, जिसे \omterm{def:g3:first-electric-circuit:battery}{बैटरी} पर लादा जाता है और
लैंपों पर \omterm{def:g1:light-and-shadows:source}{प्रकाश} तथा गरमी के रूप में उतार दिया जाता है। ख़त्म होता है
रिज़ॉर्ट के स्की करने वालों का भंडार --- यानी \omterm{def:g3:first-electric-circuit:battery}{बैटरी} की संचित \omterm{def:g5:everyday-energy:energy}{ऊर्जा}।
\end{solution}

\begin{solution}{exo:g7:circuits-series-parallel:9}
क से गुज़रने वाला हर फंदा पहले \omterm{ex:g3:first-electric-circuit:switch}{स्विच} से और फिर या तो ख या ग से गुज़रता
है --- इसलिए क के श्रेणी वाले साथी हैं \omterm{ex:g3:first-electric-circuit:switch}{स्विच} और \omterm{def:g5:series-parallel-circuits:parallel}{शाखाओं} की \emph{पूरी
जोड़ी}, पर अकेला ख या अकेला ग नहीं। ग वाली \omterm{def:g5:series-parallel-circuits:parallel}{शाखा} काट दो तो पूरी क़तार ख
में से जाती है: क जलता रहता है (ज़रा अलग ढंग से --- दो सड़कों की जगह
एक --- जैसा अगले साल के उपकरण मापेंगे)।
\end{solution}

\begin{solution}{exo:g7:circuits-series-parallel:10}
आपूर्ति के ठीक बाद साझा सड़क पर एक मुख्य \omterm{ex:g3:first-electric-circuit:switch}{स्विच}, और फिर दस समांतर
शाखाएँ, हर एक पर एक \omterm{def:g3:first-electric-circuit:bulb}{बल्ब}। समांतर का नियम हर \omterm{def:g3:first-electric-circuit:bulb}{बल्ब} को पूरी और आज़ाद चमक
देता है; और मुख्य \omterm{ex:g3:first-electric-circuit:switch}{स्विच} का श्रेणी वाला ठिकाना --- यानी वह सड़क जो हर
फंदा साझा करता है --- एक ही \omterm{ex:g3:first-electric-circuit:switch}{स्विच} को दसों पर हुक्म चलाने देता है।
\end{solution}

\begin{solution}{exo:g7:circuits-series-parallel:11}
\omterm{def:g5:series-parallel-circuits:parallel}{संधि} पर होने वाला बँटवारा बराबर-बराबर नहीं होता: किसी \omterm{def:g5:series-parallel-circuits:parallel}{शाखा} की अपनी
सड़क क़तार का विरोध कर सकती है, और लंबा पतला तार अपनी \omterm{def:g5:series-parallel-circuits:parallel}{शाखा} का हिस्सा
गला देता है --- हर सेकंड कम चलने वाले, यानी कमज़ोर चमक। सड़कें ख़ुद
धारा से लड़ती हैं: यही अगले साल के प्रतिरोध का विचार है।
\end{solution}

\begin{solution}{exo:g7:circuits-series-parallel:12}
क्ष अपनी \omterm{def:g5:series-parallel-circuits:parallel}{शाखा} पर अकेला है और उसे पूरा धक्का मिलता है: सबसे चमकीला।
त्र और ज्ञ अपनी \omterm{def:g5:series-parallel-circuits:parallel}{शाखा} का धक्का \omterm{def:g4:series-circuits:series}{श्रेणी में} बाँटते हैं: दोनों बराबर मंद,
और क्ष से नीचे। ज्ञ फुँक जाए तो उसका फंदा टूट जाता है: त्र भी उसके साथ
बुझ जाता है --- और क्ष बेपरवाह जलता रहता है। इसके बजाय क्ष फुँके तो
सिर्फ़ क्ष वाली \omterm{def:g5:series-parallel-circuits:parallel}{शाखा} मरती है: त्र और ज्ञ पहले की तरह मंद जलते रहते हैं।
\end{solution}

\begin{solution}{pb:g7:circuits-series-parallel:1}
\textbf{1.} चार समांतर शाखाएँ, हर एक पर एक ज़िम्मेदारी: \omterm{def:g5:series-parallel-circuits:parallel}{समांतर में} हर
\omterm{def:g5:series-parallel-circuits:parallel}{शाखा} पूरी ताक़त से काम करती है, बाक़ी चाहे जो करें। \omterm{def:g4:series-circuits:series}{श्रेणी में} कोई भी
जोड़ी पढ़ने वाली बत्ती के \omterm{ex:g3:first-electric-circuit:switch}{स्विच} या फुँके हुए सीढ़ी-बल्ब को विशाल लैंप
मंद या अँधेरा करने देती --- यह माफ़ी लायक नहीं।
\textbf{2.} भंडार के ठीक बाद साझा सड़क पर मुख्य \omterm{ex:g3:first-electric-circuit:switch}{स्विच}; और चारों
शाखा-स्विच अपनी-अपनी \omterm{def:g5:series-parallel-circuits:parallel}{शाखा} के भीतर, अपनी ज़िम्मेदारी के साथ \omterm{def:g4:series-circuits:series}{श्रेणी में}।
\textbf{3.} हाँ --- समांतर का नियम: किसी एक \omterm{def:g5:series-parallel-circuits:parallel}{शाखा} की दरार सिर्फ़ उसी
\omterm{def:g5:series-parallel-circuits:parallel}{शाखा} का हिस्सा रोकती है; बाक़ी की क़तार चलती रहती है।
\textbf{4.} भंडार, मुख्य \omterm{ex:g3:first-electric-circuit:switch}{स्विच}, फिर एक \omterm{def:g5:series-parallel-circuits:parallel}{संधि} से चार शाखाएँ ---
(विशाल लैंप $+$ \omterm{ex:g3:first-electric-circuit:switch}{स्विच}), (कोहरे का भोंपू $+$ \omterm{ex:g3:first-electric-circuit:switch}{स्विच}), (पढ़ने वाली बत्ती
$+$ \omterm{ex:g3:first-electric-circuit:switch}{स्विच}), (सीढ़ियों की बत्ती $+$ \omterm{ex:g3:first-electric-circuit:switch}{स्विच}) --- जो मिलकर भंडार तक लौट
जाती हैं।
\textbf{5.} हाँ: हर \omterm{def:g5:series-parallel-circuits:parallel}{शाखा} भंडार के दोनों पक्षों के बीच फैली है, इसलिए
हर एक को पूरा धक्का मिलता है, चाहे साथ में कितने ही पड़ोसी पी रहे हों।
\textbf{6.} नहीं: विशाल लैंप में से गुज़रती धारा उसे पूरी छोड़कर जाती
है --- कुर्सियाँ बिना खाए चक्कर लगाती हैं। लैंप जो लेता है वह भंडार की
\omterm{def:g5:everyday-energy:energy}{ऊर्जा} है, जिसे गुज़रती क़तार पहुँचा देती है।
\textbf{7.} एक साथ चार पूरी क़तारें: भंडार की \omterm{def:g5:everyday-energy:energy}{ऊर्जा} चार ज़िम्मेदारियों
में तेज़ी से बहती है। बिल \omterm{def:g3:first-electric-circuit:battery}{बैटरी} वाले कमरे में चुकाया जाता है ---
संचित \omterm{def:g5:everyday-energy:energy}{ऊर्जा} में, कभी खोई हुई धारा में नहीं।
\textbf{8.} भंडार के $+$ से निकलकर, मुख्य \omterm{ex:g3:first-electric-circuit:switch}{स्विच} से होकर, \omterm{def:g5:series-parallel-circuits:parallel}{संधि} तक; फिर
पढ़ने वाली बत्ती की \omterm{def:g5:series-parallel-circuits:parallel}{शाखा} में, उसके बंद \omterm{ex:g3:first-electric-circuit:switch}{स्विच} और लैंप से होकर; उस \omterm{def:g5:series-parallel-circuits:parallel}{संधि}
से आगे जहाँ बाक़ी तीनों (अँधेरी) शाखाएँ फिर मिलती हैं; और भंडार के $-$
पर वापस। वह विशाल लैंप, भोंपू और सीढ़ी वाली \omterm{def:g5:series-parallel-circuits:parallel}{शाखाओं} से पूरी तरह बचकर
निकल जाती है।
\textbf{9.} और कुछ नहीं --- सिर्फ़ सीढ़ी वाली \omterm{def:g5:series-parallel-circuits:parallel}{शाखा} अपने \omterm{def:g3:first-electric-circuit:bulb}{बल्ब} के साथ
मरती है: यही समांतर के नियम की स्वतंत्रता है।
\textbf{10.} हर \omterm{def:g5:series-parallel-circuits:parallel}{शाखा} से पहले लगा एक \omterm{def:g6:circuits-and-safety:short}{लघुपथन}। क़तार हर काम की सड़क छोड़कर
उस बिना-मेहनत वाले नंगे शॉर्टकट पर चली जाती है: लैंप और भोंपू भूखे रह
जाते हैं जबकि शॉर्टकट वाला तार बेतहाशा, तपती दौड़ ढोता है --- और तार
हटने तक हर ज़िम्मेदारी एक साथ नाकाम रहती है।
\textbf{11.} विशाल लैंप को दूसरे भंडार पर उसका अपना स्वतंत्र फंदा दे
दो: भंडार दो, उसका अपना \omterm{ex:g3:first-electric-circuit:switch}{स्विच}, और विशाल लैंप --- और आम पटल के साथ कोई
तार साझा नहीं। तब भंडार एक पर लगा शॉर्टकट सब कुछ चुप कर देगा,
\emph{सिवाय} जहाज़ों वाली बत्ती के।
\textbf{12.} जैसे: ``आज रात हर ज़िम्मेदारी को समांतर के नियम ने अपना
राजा बनाए रखा। रात ने एक बार फिर उस भरोसे का खंडन किया कि लैंप धारा
निगल जाते हैं --- भंडार चुकते रहे, फिर भी क़तार पूरी लौटती रही। और एक
नंगे तार ने, भंडार को शॉर्ट करके, हमें हमारा सबसे बुरा घंटा दिया: हर
सड़क ख़ाली, और शॉर्टकट जलता हुआ।''
\end{solution}
